\section{関連研究}\label{sec:related}

\subsection{自己改善と共進化}\label{sec:related:treesearch}

自己反省や自己洗練の研究では、モデルが自ら生成した批評を後続の試行へ
戻す \cite{madaan2023self,shinn2023reflexion}。これらは成果物や推論過程
を改善するが、評価者と権限構造はおおむね固定されている。Constitutional
ARI-RQGM は、指示文で定義された役割と効用規則も版管理された制度上の
成果物として扱い、その置換に証拠、役割分離、適法な遷移を求める。

最も近い概念的な祖先は Red Queen G\"odel Machine
\cite{iacob2026red} である。この方式では、主体と評価者が共進化し、
採点関数自体も変化する。ARI-RQGM はこの共進化と敵対的な再試験を
取り入れ、さらに固定憲法、期ごとの不可分な更新、台帳の単一書き手、
責任対象を明示した制裁、隔離再生成、可変の効用規則と不変の検証処理
を組み合わせる。元の方式に既にある共進化と期境界更新を本稿の新規性
とは数えない。本稿の概念的な差分は、変更可能な制度を権限表と構成要素
状態遷移表で囲み、退役後の影響修復と固定証拠判定まで同じ来歴上で結ぶ
ことである。個々の部品ではなく、この統合とその検査可能な境界が主な
システム上の貢献である。

Constitutional AI は、文章化した原則を用いる自己批評と人工知能からの
帰還を示した \cite{bai2022constitutional}。本稿の「憲法」は、モデル
出力の望ましさを教える規範だけではなく、モデルが迂回できない権限表、
状態遷移、証拠形式を固定処理で執行する点が異なる。AI Control は、意図
して監視を逃れようとする非信頼モデルを前提に、信頼モデルとの分離や
監視手続きを評価する \cite{greenblatt2023ai}。本稿はそれより狭く、
共謀に対する定量的な制御評価をまだ行っていないが、執筆、査読、裁定、
台帳更新を同一権限へ集めない設計を採る。

\subsection{自律研究システム}\label{sec:related:autonomous}

AI Scientist は着想から実験、原稿作成までを一貫して自動化する流れを
示し \cite{lu2024ai}、後続研究は自動科学発見へ向かう広い展望を整理
した \cite{lu2026towards}。AIDE は実験コードの改善を反復的な主体の
問題として扱う \cite{jiang2025aide}。これらは、永続的な成果物、実行
可能な評価、自動執筆の価値を示す。主な最適化対象は研究成果物である。
これに対し ARI-RQGM は、その成果物を作り評価する制度を対象とする。

Agent Laboratory は研究助手として協働する複数のモデル役を調べ
\cite{schmidgall2025agent}、MLE-bench は機械学習開発を自動化する主体
向けの実行可能な課題を提供する \cite{chan2024mle}。下位の研究実行系
が連鎖的な推論促進法 \cite{wei2022chain} を使うか、大規模な候補生成
\cite{li2022competition} を使うかは、本方式の憲法境界に影響しない。
統治対象は内部の生成手順ではなく、稼働中の役割、その権限、判断根拠
である。

Story2Proposal は、物語形式の科学的主張を実行可能な提案と証拠へ
つなぐ契約を論じる \cite{qian2026story2proposal}。統治された論文
保管庫も同じ原理を使い、実行証拠を執筆役から独立した固定基準とする。
査読役が進化しても、実行結果に反する記述を正しいことにはできない。

\subsection{評価と証拠}\label{sec:related:eval}

PaperBench は構造化された採点基準を使い、主体が研究論文を再現できる
かを評価する \cite{starace2025paperbench}。ここから得られる重要な教訓
は、評価上の主張には明確な課題構造、観測可能な成果物、独立した採点
境界が必要だという点である。Constitutional ARI-RQGM はこの考えを
最終成果だけでなく制度史へ広げ、攻撃、弁護、裁定、遷移、修復を来歴
付きの型の定まった記録として残す。

敵対的評価は、固定されたモデルや方策を強化する目的で使われることが
多い。本方式では、検証済み攻撃に二つの役割を持たせる。第一に稼働中
の構成要素へ責任を結び付け、第二に後継候補を試す過去事例とする。
同じ故障事例が、現職の問題を示す証拠であると同時に、後継の採用試験
にもなる。

\subsection{本研究の位置づけ}\label{sec:related:foundations}

最小権限、完全仲介、権限分離は古典的な保護設計原則である
\cite{saltzer1975protection}。決定論的な状態機械は、確定した入力列から
同じ状態を再構築する障害回復の基礎を与える
\cite{schneider1990implementing}。追記記録と決定論的な射影を主体の状態
管理へ用いる近年の事象駆動設計も同じ系譜にある
\cite{nakajima2026the}。したがって、単一書き手、状態再生、追記履歴を
それ自体の新発明とは主張しない。技術上の課題は、これらを共進化する
RQGM の期、責任帰属、退役修復、論文証拠判定へ矛盾なく接続することに
ある。

\href{https://www.w3.org/TR/prov-o/}{W3C PROV-O 勧告}は、実体・活動・
主体からなる相互運用可能な来歴を定める。
\href{https://www.usenix.org/conference/usenixsecurity19/presentation/torres-arias}{in-toto}
は供給工程の署名付き証拠を扱い、
\href{https://docs.sigstore.dev/logging/overview/}{Sigstore Rekor} は第三者
が監査できる追記型透明性記録を提供する。これらには、現版の局所的な未署名
保存点にない相互運用性や暗号学的・外部監査上の強みがある。一方、RQGM
固有の退役影響閉包や効用再採点は定めない。

\begin{table*}[t]
  \centering
  \scriptsize
  \begin{tabularx}{\textwidth}{@{}p{0.12\textwidth}*{7}{X}@{}}
    \toprule
    \textbf{研究系譜} & \textbf{更新対象} & \textbf{固定基準} &
    \textbf{役割権限} & \textbf{不可分な状態更新} &
    \textbf{履歴完全性} & \textbf{修復} & \textbf{脅威・証拠} \\
    \midrule
    RQGM \cite{iacob2026red} & 主体、評価者、効用 &
      期内評価 & 主眼ではない & 期ごとの更新 &
      主眼ではない & 共進化選択 &
      概念・実験上の共進化 \\
    事象駆動主体 \cite{nakajima2026the} & グラフの挙動と状態 &
      出来事契約 & 意図と状態の分離 & 決定論的射影 &
      追記来歴 & 分岐・再生 &
      監査可能性の設計 \\
    Story2Proposal \cite{qian2026story2proposal} & 科学記述と提案 &
      構造化契約 & 主眼ではない & 契約段階 &
      証拠の結合 & 契約による改稿 &
      執筆支援の評価 \\
    AI Control \cite{greenblatt2023ai} & 制御手続き &
      信頼監視・手続き & 信頼・非信頼の分離 &
      手続き依存 & 監視記録 & 手続き上の応答 &
      意図的回避と実証制御 \\
    来歴・透明性方式 & 実体、活動、署名済み記述 &
      来歴形式、鍵、記録手順 & 認証された生成者・検証者 &
      署名工程または追記 & 相互運用または外部監査可能な履歴 &
      検証（意味的退役修復ではない） & 暗号学的完全性 \\
    本方式 & 指示文、構成要素、効用規則 &
      固定カーネル・検証器 & 処理内の権限行列 &
      境界取引 & 意味項目を結ぶハッシュ値、外部固定なし &
      退役閉包と再採点 &
      五故障例と回帰試験、共謀実験なし \\
    \bottomrule
  \end{tabularx}
  \caption{構成要素単位の比較。記載があることは、目的の異なる方式を
    直接比較できる証明ではない。本稿の新規性は個別機構の発明ではなく、
    統合に置く。}
  \label{tab:related-comparison}
\end{table*}

透明性履歴とソフトウェア供給網の来歴は、どのバイトと身元を固定したかを
問い、参照監視器と能力方式は、どの主体がどの対象へ何を行えるかを問う。
安全性を監視する状態機械は、すべての実行列が安全言語へ属するかを
問う \cite{schneider2000enforceable}。これらは本稿が依拠する系譜であり、新規性ではない。現実装は最小
権限と完全仲介 \cite{saltzer1975protection}、決定論的な状態射影
\cite{schneider1990implementing} を取り入れるが、独立監査者に固定した
履歴先頭値も OS で強制した役割能力もない。また、期、取引、動議枠、修復
を合成した状態は形式的モデル検査を受けていない。

Constitutional ARI-RQGM は、自律科学一般での優越性を主張するものでも、
自己改変する主体の問題をすべて解くものでもない。扱う問いはより狭い。
モデルで定義した研究役割と効用規則を実行中に変更しながら、役割自身
には採点規則の免除、失敗履歴の削除、後継の直接任命を許さない制度を
どう構成するか、という問いである。貢献は、共進化、決定論的な憲法
執行、不可分な状態更新、来歴を考慮した修復、固定された証拠判定を
一つの運用可能な仕組みにまとめた点にある。
